% Auto-generated from raw.txt
% Usage:
%   \vocabentry{word}{/ipa/ (pos)}{definition}{example}{note}

\vocabentry{Innovation}
{/ˌɪnəˈveɪʃn/ (n)}
{The action or process of innovating; a new method, idea, or product.}
{Continuous innovation in the energy sector is our best hope for a cleaner and more prosperous world.}
{Đây là một danh từ dùng để chỉ sự đổi mới sáng tạo, chủ đề cốt lõi của ngành năng lượng}

\vocabentry{Renewable}
{/rɪˈnuːəbl/ (adj)}
{(Of a natural resource or source of energy) not depleted when used.}
{Countries are investing heavily in renewable energy sources like wind and solar to combat climate change.}
{Đây là một tính từ dùng để chỉ năng lượng tái tạo, có khả năng phục hồi tự nhiên.}

\vocabentry{Photovoltaic}
{/ˌfoʊtoʊvɔːlˈteɪɪk/ (adj)}
{Relating to the production of electric current at the junction of two substances exposed to light.}
{Photovoltaic panels convert sunlight directly into electricity for residential use.}
{Đây là một tính từ dùng để chỉ công nghệ điện quang/quang điện, thường dùng cho các tấm pin mặt trời.}

\vocabentry{Geothermal}
{/ˌdʒiːoʊˈθɜːrml/ (adj)}
{Relating to or produced by the internal heat of the earth.}
{Iceland is a world leader in utilizing geothermal energy for heating and power.}
{Đây là một tính từ dùng để chỉ địa nhiệt, năng lượng từ lòng đất.}

\vocabentry{Hydrogen fuel cell}
{/ˈhaɪdrədʒən ˈfjuːəl sel/ (n)}
{A device that converts chemical energy from a fuel into electricity through a chemical reaction of hydrogen with oxygen.}
{Hydrogen fuel cells are being developed as a clean alternative for powering heavy trucks and ships.}
{Đây là một danh từ dùng để chỉ pin nhiên liệu hydro.}

\vocabentry{Decarbonization}
{/ˌdiːˌkɑːrbənaɪˈzeɪʃn/ (n)}
{The reduction or removal of carbon dioxide emissions from a process or economy.}
{The global shift towards decarbonization is essential for achieving net-zero targets.}
{Đây là một danh từ dùng để chỉ quá trình giảm phát thải carbon.}

\vocabentry{Efficiency}
{/ɪˈfɪʃnsi/ (n)}
{The state or quality of being efficient; achieving maximum productivity with minimum wasted effort or expense.}
{New power electronics are designed for maximum efficiency and minimum heat loss.}
{Đây là một danh từ dùng để chỉ hiệu suất hoặc sự tối ưu hóa tài nguyên năng lượng}

\vocabentry{Grid}
{/ɡrɪd/ (n)}
{A network of lines that cross each other to form a series of squares or rectangles; in energy, the system of power lines and stations.}
{Modernizing the aging power grid is necessary to integrate intermittent renewable energy.}
{Đây là một danh từ dùng để chỉ mạng lưới điện quốc gia hoặc khu vực.}

\vocabentry{Turbine}
{/ˈtɜːrbaɪn/ (n)}
{A machine for producing continuous power in which a wheel or rotor is made to revolve by a fast-moving flow of water, steam, gas, or air.}
{Wind turbines are becoming more efficient at capturing energy even at low wind speeds.}
{Đây là một danh từ dùng để chỉ tua-bin.}

\vocabentry{Hydropower}
{/ˈhaɪdroʊpaʊər/ (n)}
{Electricity produced from the energy of moving water.}
{Hydropower remains one of the largest sources of low-carbon electricity globally.}
{Đây là một danh từ dùng để chỉ thủy điện.}

\vocabentry{Nuclear fusion}
{/ˈnuːkliər ˈfjuːʒn/ (n)}
{A nuclear reaction in which atomic nuclei of low atomic number fuse to form a heavier nucleus with the release of energy.}
{Scientists hope that nuclear fusion will eventually provide a limitless source of clean energy.}
{Đây là một cụm danh từ dùng để chỉ phản ứng nhiệt hạch, hy vọng cho năng lượng sạch vô hạn.}

\vocabentry{Biofuel}
{/ˈbaɪoʊfjuːəl/ (n)}
{A fuel derived immediately from living matter.}
{Biofuel produced from algae is a promising alternative to traditional jet fuels.}
{Đây là một danh từ dùng để chỉ nhiên liệu sinh học làm từ thực vật hoặc rác thải hữu cơ.}

\vocabentry{Storage}
{/ˈstɔːrɪdʒ/ (n)}
{The action or method of storing something for future use.}
{Energy storage systems are vital for stabilizing the supply of solar power during the night.}
{Trong năng lượng, đây là việc lưu trữ năng lượng, ví dụ bằng pin.}

\vocabentry{Sustainable}
{/səˈsteɪnəbl/ (adj)}
{Able to be maintained at a certain rate or level.}
{We must transition to a sustainable energy system to protect the planet for future generations.}
{Đây là một tính từ dùng để chỉ sự bền vững về năng lượng và môi trường.}

\vocabentry{Emission}
{/iˈmɪʃn/ (n)}
{The production and discharge of something, especially gas or radiation.}
{Reducing methane emissions is a top priority in energy innovation.}
{Đây là một danh từ dùng để chỉ khí thải, thường là CO2.}

\vocabentry{Infrastructure}
{/ˈɪnfrəstrʌktʃər/ (n)}
{The basic physical and organizational structures and facilities needed for the operation of a society or enterprise.}
{Massive investment in infrastructure is required to support the transition to green energy.}
{Đây là một danh từ dùng để chỉ hệ thống hạ tầng kỹ thuật, bao gồm lưới điện và các công trình năng lượng}

\vocabentry{Solar farm}
{/ˈsoʊlər fɑːrm/ (n)}
{An area of land on which a large number of solar panels are deployed to generate electricity.}
{Large solar farms in the desert can provide electricity for thousands of homes.}
{Đây là một cụm danh từ dùng để chỉ cánh đồng năng lượng mặt trời.}

\vocabentry{Offshore wind}
{/ˌɔːfʃɔːr ˈwɪnd/ (n)}
{Wind power generated from wind turbines built over bodies of water, usually the ocean.}
{Offshore wind has the advantage of stronger and more consistent winds compared to onshore sites.}
{Đây là một cụm danh từ dùng để chỉ điện gió ngoài khơi.}

\vocabentry{Smart meter}
{/ˌsmɑːrt ˈmiːtər/ (n)}
{An electronic device that records consumption of electric energy and communicates that information for monitoring and billing.}
{Smart meters allow consumers to track their energy usage in real-time and reduce waste.}
{Đây là một cụm danh từ dùng để chỉ công tơ điện thông minh.}

\vocabentry{Carbon capture}
{/ˈkɑːrbən ˈkæptʃər/ (n)}
{The process of trapping waste carbon dioxide at its source and depositing it where it will not enter the atmosphere.}
{Carbon capture technology is being tested at coal-fired power plants to mitigate climate impact.}
{Đây là một cụm danh từ dùng để chỉ công nghệ thu giữ carbon.}

\vocabentry{Intermittent}
{/ˌɪntərˈmɪtənt/ (adj)}
{Occurring at irregular intervals; not continuous or steady.}
{Intermittent energy sources need to be complemented by stable base-load power.}
{Đây là một tính từ dùng để chỉ tính chất gian đoạn, không liên tục của điện gió và điện mặt trời}

\vocabentry{Silicon}
{/ˈsɪlɪkən/ (n)}
{A chemical element used in the manufacture of solar cells and computer chips.}
{Advancements in silicon technology have dramatically reduced the cost of solar energy.}
{Đây là một danh từ dùng để chỉ silic, vật liệu chủ đạo trong pin mặt trời.}

\vocabentry{Semi-conductor}
{/ˌsemikənˈdʌktər/ (n)}
{A solid substance that has a conductivity between that of an insulator and that of most metals.}
{High-efficiency semi-conductors are essential for converting DC power from solar panels to AC power.}
{Đây là một danh từ dùng để chỉ chất bán dẫn, cốt lõi của thiết bị điện tử năng lượng.}

\vocabentry{Sustainability}
{/səˌsteɪnəˈbɪləti/ (n)}
{The ability to be maintained at a certain rate or level; avoidance of the depletion of natural resources.}
{The sustainability of renewable energy projects depends on responsible mineral sourcing.}
{Đây là một danh từ dùng để chỉ tính bền vững, khả năng duy trì phát triển lâu dài mà không cạn kiệt tài nguyên}

\vocabentry{Tidal power}
{/ˈtaɪdl paʊər/ (n)}
{A form of hydropower that converts the energy obtained from tides into useful forms of power.}
{Tidal power is highly predictable, making it a reliable source of renewable electricity.}
{Đây là một cụm danh từ dùng để chỉ năng lượng thủy triều.}

\vocabentry{Thermal}
{/ˈθɜːrml/ (adj)}
{Relating to heat.}
{Thermal energy can be stored in molten salts to generate electricity after sunset.}
{Đây là một tính từ dùng để chỉ các yếu tố thuộc về nhiệt năng.}

\vocabentry{Electrification}
{/ɪˌlektrɪfɪˈkeɪʃn/ (n)}
{The action of process of charging something with electricity. (Đây là một danh từ dùng để chỉ quá trình điện khí hóa (ví dụ: chuyển từ xe xăng sang xe điện).).}
{The electrification of public transport is a key step towards cleaner cities.}
{}

\vocabentry{Microgrid}
{/ˈmaɪkroʊɡrɪd/ (n)}
{A small network of electricity users with a local source of supply that is usually attached to a centralized national grid.}
{Microgrids can provide reliable power to remote communities or during natural disasters.}
{Đây là một danh từ dùng để chỉ mạng lưới điện siêu nhỏ cục bộ.}

\vocabentry{Battery pack}
{/ˈbætəri pæk/ (n)}
{A set of any number of identical batteries or individual battery cells.}
{The battery pack in an electric car determines its driving range.}
{Đây là một cụm danh từ dùng để chỉ bộ pin năng lượng.}

\vocabentry{Megawatt}
{/ˈmeɡəwɑːt/ (n)}
{A unit of power equal to one million watts.}
{The new wind farm has a total capacity of five hundred megawatts.}
{Đây là một danh từ dùng để chỉ đơn vị đo công suất điện quy mô lớn.}

\vocabentry{Fossil fuel}
{/ˈfɑːsl fjuːəl/ (n)}
{A natural fuel such as coal or gas, formed in the geological past from the remains of living organisms.}
{We must reduce our dependency on fossil fuels to meet global environmental goals.}
{Đây là một cụm danh từ dùng để chỉ nhiên liệu hóa thạch.}

\vocabentry{Transition}
{/trænˈzɪʃn/ (n)}
{The process or a period of changing from one state or condition to another.}
{The transition period between coal and gas will vary from country to country.}
{Đây là một danh từ dùng để chỉ sự chuyển đổi giai đoạn, đặc biệt từ nhiên liệu hóa thạch sang năng lượng sạch}

\vocabentry{Biomass}
{/ˈbaɪoʊmæs/ (n)}
{Organic matter used as a fuel, especially in a power station.}
{Biomass energy utilizes wood, agricultural waste, and even garbage to produce heat.}
{Đây là một danh từ dùng để chỉ sinh khối, vật liệu hữu cơ dùng làm năng lượng.}

\vocabentry{Incentive}
{/ɪnˈsentɪv/ (n)}
{A thing that motivates or encourages someone to do something.}
{Tax incentives for solar panel installation have encouraged many homeowners to go green.}
{Trong năng lượng là các chính sách khuyến khích từ chính phủ.}

\vocabentry{Optimization}
{/ˌɑːptɪməˈzeɪʃn/ (n)}
{The action of making the best or most effective use of a situation or resource.}
{Artificial intelligence is playing a huge role in the optimization of the modern energy grid.}
{Đây là một danh từ dùng để chỉ quá trình tối ưu hóa hiệu suất hệ thống năng lượng}

\vocabentry{Interconnected}
{/ˌɪntərkəˈnektɪd/ (adj)}
{(Of two or more things) connected with each other.}
{An interconnected market allows for more stable electricity prices across regions.}
{Đây là một tính từ dùng để chễ tính liên kết với nhau, đặc biệt là các hệ thống lưới điện khu vực}

\vocabentry{Photons}
{/ˈfoʊtɑːnz/ (n)}
{Particles representing a quantum of light or other electromagnetic radiation.}
{Solar cells work by absorbing photons from sunlight to release electrons.}
{Đây là một danh từ dùng để chỉ các hạt ánh sáng trong vật lý năng lượng mặt trời.}

\vocabentry{Transmission}
{/trænzˈmɪʃn/ (n)}
{The action or process of transmitting something.}
{High-voltage transmission lines are needed to move electricity from wind farms to cities.}
{Trong năng lượng là sự truyền tải điện năng đi xa.}

\vocabentry{Distribution}
{/ˌdɪstrɪˈbjuːʃn/ (n)}
{The action of sharing something out among a number of recipients.}
{Smart distribution networks can automatically detect and isolate power outages.}
{Trong năng lượng là sự phân phối điện đến từng hộ gia đình.}

\vocabentry{Retrofit}
{/ˈretroʊfɪt/ (v)}
{Add (a component or accessory) to something that did not have it when manufactured.}
{The old factory was retrofitted with energy-efficient lighting and insulation.}
{Đây là một động từ dùng để chỉ việc cải tạo thiết bị cũ bằng công nghệ năng lượng mới.}

\vocabentry{Insulation}
{/ˌɪnsəˈleɪʃn/ (n)}
{Material used to prevent the loss of heat or the intrusion of sound.}
{Proper roof insulation can cut heating bills by up to thirty percent.}
{Đây là một danh từ dùng để chỉ vật liệu cách nhiệt, giúp tiết kiệm năng lượng.}

\vocabentry{Renewability}
{/rɪˈnuːəˈbɪlɪti/ (n)}
{The quality of being capable of being renewed.}
{The renewability of wood depends on responsible forestry practices.}
{Đây là một danh từ dùng để chỉ khả năng tái tạo.}

\vocabentry{Conduction}
{/kənˈdʌkʃn/ (n)}
{The process by which heat or electricity is directly transmitted through a substance.}
{Copper is widely used in energy systems due to its excellent electrical conduction.}
{Đây là một danh từ dùng để chỉ sự dẫn điện hoặc dẫn nhiệt.}

\vocabentry{Generator}
{/ˈdʒenəreɪtər/ (n)}
{A machine that converts mechanical energy into electricity.}
{During the power cut, the hospital relied on its backup diesel generator.}
{Đây là một danh từ dùng để chỉ máy phát điện.}

\vocabentry{Reactor}
{/riˈæktər/ (n)}
{A device or complex in which a nuclear reaction is maintained and controlled.}
{Modern nuclear reactors are designed with multiple safety systems to prevent leaks.}
{Đây là một danh từ dùng để chỉ lò phản ứng hạt nhân.}

\vocabentry{Prototype}
{/ˈproʊtətaɪp/ (n)}
{A first, typical or preliminary model of something.}
{The engineers are testing a prototype of a new high-capacity battery.}
{Đây là một danh từ dùng để chỉ mẫu thử nghiệm của một thiết bị năng lượng mới.}

\vocabentry{Storage capacity}
{/ˈstɔːrɪdʒ kəˈpæsəti/ (n)}
{The amount of energy that can be stored in a system.}
{Increasing the storage capacity of the grid is essential for a 100\% renewable future.}
{Đây là một cụm danh từ dùng để chỉ dung lượng lưu trữ.}

\vocabentry{Breakthrough}
{/ˈbreɪkθruː/ (n)}
{A sudden, dramatic, and important discovery or development.}
{A breakthrough in solid-state batteries could make electric cars much lighter and safer.}
{Đây là một danh từ dùng để chỉ sự đột phá quan trọng trong nghiên cứu và công nghệ năng lượng}

\vocabentry{Carbon footprint}
{/ˌkɑːrbən ˈfʊtprɪnt/ (n)}
{The amount of carbon dioxide and other carbon compounds emitted due to the consumption of fossil fuels by a person or group.}
{Switching to renewable energy is the fastest way to reduce your carbon footprint.}
{Đây là một cụm danh từ dùng để chỉ dấu chân carbon.}

\vocabentry{Net-zero}
{/ˌnet ˈzɪroʊ/ (adj)}
{Achieving a balance between the greenhouse gases put into the atmosphere and those taken out.}
{Many corporations have pledged to become net-zero by the year 2050.}
{Đây là một tính từ dùng để chỉ mục tiêu phát thải ròng bằng không.}

\vocabentry{Adaptation}
{/ˌædæpˈteɪʃn/ (n)}
{The action or process of adapting or being adapted.}
{Adaptation to rising sea levels is already a priority for many coastal power plants.}
{Đây là một danh từ dùng để chỉ sự thích nghi của con người và hệ thống với biến đổi khí hậu và môi trường}

\vocabentry{Feasibility}
{/ˌfiːzəˈbɪləti/ (n)}
{The state or degree of being easily or conveniently done.}
{The feasibility of using hydrogen for home heating is still being debated by experts.}
{Đây là một danh từ dùng để chỉ tính khả thi, sự thực thi thực tế của một dự án năng lượng}

\vocabentry{Regulation}
{/ˌreɡjuˈleɪʃn/ (n)}
{A rule or directive made and maintained by an authority.}
{Environmental regulations have forced many coal plants to shut down or upgrade.}
{Đây là một danh từ dùng để chỉ quy chuẩn, quy định hoặc luật lệ về năng lượng và môi trường}

\vocabentry{Subsidize}
{/ˈsʌbsɪdaɪz/ (v)}
{Support (an organization or activity) financially.}
{Many governments subsidize the purchase of electric cars to reduce pollution.}
{Đây là một động từ dùng để chỉ việc trợ cấp từ chính phủ cho năng lượng sạch.}

\vocabentry{Thermodynamics}
{/ˌθɜːrmoʊdaɪˈnæmɪks/ (n)}
{The branch of physical science that deals with the relations between heat and other forms of energy.}
{Thermodynamics explains how heat can be converted into mechanical work in an engine.}
{Đây là một danh từ dùng để chỉ nhiệt động lực học.}

\vocabentry{Kinetic}
{/kɪˈnetɪk/ (adj)}
{Relating to or resulting from motion.}
{A wind turbine converts the kinetic energy of the wind into electrical energy.}
{Đây là một tính từ dùng để chỉ động năng, ví dụ từ gió hoặc nước chảy.}

\vocabentry{Conservation}
{/ˌkɑːnsərˈveɪʃn/ (n)}
{The action of conserving something, in particular preservation, protection, or restoration of the natural environment.}
{Water conservation is becoming increasingly important for hydropower production during droughts.}
{Đây là một danh từ dùng để chỉ sự bảo tồn, tiết kiệm tài nguyên năng lượng}

\vocabentry{Off-grid}
{/ˌɔːf ˈɡrɪd/ (adj)}
{Not using or depending on public utilities, especially the supply of electricity.}
{Solar panels allow many remote cabins to operate entirely off-grid.}
{Đây là một tính từ dùng để chỉ trạng thái tự cấp năng lượng, không phụ thuộc lưới điện.}

\vocabentry{Lithium-ion}
{/ˌlɪθiəm ˈaɪən/ (n)}
{A type of rechargeable battery in which lithium ions move from the negative electrode to the positive electrode during discharge.}
{Lithium-ion batteries are currently the standard for smartphones and electric vehicles.}
{Đây là một danh từ dùng để chỉ pin Lithium-ion phổ biến hiện nay.}

\vocabentry{Electrolysis}
{/ɪˌlekˈtrɑːləsɪs/ (n)}
{Chemical decomposition produced by passing an electric current through a liquid or solution containing ions.}
{Green hydrogen is produced through the electrolysis of water using renewable electricity.}
{Đây là một danh từ dùng để chỉ quá trình điện phân nước để sản xuất hydro.}

\vocabentry{Fission}
{/ˈfɪʃn/ (n)}
{The action of dividing or splitting something into two or more parts; specifically nuclear fission.}
{Nuclear fission provides a steady supply of carbon-free electricity in many countries.}
{Đây là một danh từ dùng để chỉ phản ứng phân hạch hạt nhân.}

\vocabentry{Geopolitical}
{/ˌdʒiːoʊpəˈlɪtɪkl/ (adj)}
{Relating to politics, especially international relations, as influenced by geographical factors.}
{The transition to renewables could shift the geopolitical balance away from oil-producing nations.}
{Đây là một tính từ dùng để chỉ các yếu tố địa chính trị ảnh hưởng đến nguồn cung năng lượng.}

\vocabentry{Scalability}
{/ˌskeɪləˈbɪləti/ (n)}
{The capacity to be changed in size or scale.}
{The scalability of offshore wind power could allow it to meet a huge part of global demand.}
{Đây là một danh từ dùng để chỉ khả năng mở rộng quy mô nhanh chóng của một giải pháp năng lượng}

\vocabentry{Interoperability}
{/ˌɪntərˌɑːpərəˈbɪləti/ (n)}
{The ability of computer systems or software to exchange and make use of information.}
{Interoperability between different electric vehicle charging networks is essential for users.}
{Trong năng lượng là khả năng tương thích giữa các thiết bị sạc và xe khác nhau.}

\vocabentry{Voltage}
{/ˈvoʊltɪdʒ/ (n)}
{An electromotive force or potential difference expressed in volts.}
{Step-up transformers are used to increase the voltage for long-distance transmission.}
{Đây là một danh từ dùng để chỉ điện áp.}

\vocabentry{Current}
{/ˈkɜːrənt/ (n)}
{A flow of electricity which results from the ordered directional movement of electrically charged particles.}
{Direct current (DC) from solar panels must be converted to alternating current (AC) for home use.}
{Đây là một danh từ dùng để chỉ dòng điện.}

\vocabentry{Transformer}
{/trænsˈfɔːrmər/ (n)}
{An apparatus for reducing or increasing the voltage of an alternating current.}
{Transformers are a critical component of the electrical distribution system.}
{Đây là một danh từ dùng để chỉ máy biến áp.}

\vocabentry{Capacitor}
{/kəˈpæsɪtər/ (n)}
{A device used to store an electric charge, consisting of one or more pairs of conductors separated by an insulator.}
{Super-capacitors can charge and discharge energy much faster than traditional batteries.}
{Đây là một danh từ dùng để chỉ tụ điện.}

\vocabentry{Resistance}
{/rɪˈzɪstəns/ (n)}
{The degree to which a substance or device opposes the passage of an electric current.}
{Choosing materials with low electrical resistance helps to minimize energy loss in cables.}
{Đây là một danh từ dùng để chỉ điện trở.}

\vocabentry{Induction}
{/ɪnˈdʌkʃn/ (n)}
{The production of an electromotive force across an electrical conductor in a changing magnetic field.}
{Wireless charging uses electromagnetic induction to transfer energy to smartphones.}
{Đây là một danh từ dùng để chỉ hiện tượng cảm ứng điện từ.}

\vocabentry{Photocell}
{/ˈfoʊtoʊsel/ (n)}
{A device that generates an electric current or voltage when exposed to light.}
{Photocells are often used in automatic street lighting to detect dusk and dawn.}
{Đây là một danh từ dùng để chỉ tế bào quang điện.}

\vocabentry{Refining}
{/rɪˈfaɪnɪŋ/ (n)}
{The process of removing impurities or unwanted elements from a substance.}
{Lithium refining is a critical step in the supply chain for battery manufacturing.}
{Trong năng lượng là quá trình lọc dầu hoặc tinh chế quặng.}

\vocabentry{Extraction}
{/ɪkˈstrækʃn/ (n)}
{The action of taking out something, especially using effort or force.}
{The extraction of rare earth metals has significant environmental consequences.}
{Đây là một danh từ dùng để chỉ sự khai thác tài nguyên như dầu mỏ hoặc khoáng sản.}

\vocabentry{Combustion}
{/kəmˈbʌstʃən/ (n)}
{The process of burning something.}
{Internal combustion engines are being phased out in favor of electric motors.}
{Đây là một danh từ dùng để chỉ sự đốt cháy nhiên liệu.}

\vocabentry{Pipeline}
{/ˈpaɪplaɪn/ (n)}
{A long pipe, typically underground, for conveying oil, gas, etc. over long distances.}
{A new pipeline was built to transport natural gas from the sea to the power plant.}
{Đây là một danh từ dùng để chỉ đường ống dẫn năng lượng.}

\vocabentry{Resilience}
{/rɪˈzɪliens/ (n)}
{The capacity to recover quickly from difficulties; toughness.}
{Building a more resilient grid is a top priority for governments worldwide.}
{Đây là một danh từ dùng để chỉ khả năng chống chịu và phục hồi của hệ thống năng lượng trước các sự cố}

\vocabentry{Volatility}
{/ˌvɑːləˈtɪləti/ (n)}
{Liability to change rapidly and unpredictably, especially for the worse.}
{The volatility of oil prices makes it difficult for companies to plan long-term budgets.}
{Đây là một danh từ dùng để chỉ sự biến động giá năng lượng.}

\vocabentry{Subsidy}
{/ˈsʌbsədi/ (n)}
{A sum of money granted by the government or a public body to assist an industry.}
{Reducing fossil fuel subsidies could accelerate the transition to green energy.}
{Đây là một danh từ dùng để chỉ tiền trợ cấp.}

\vocabentry{Diversification}
{/daɪˌvɜːrsɪfɪˈkeɪʃn/ (n)}
{The process of becoming more varied.}
{Energy diversification is essential to reduce reliance on imported natural gas.}
{Đây là một danh từ dùng để chỉ sự đa dạng hóa nguồn năng lượng để đảm bảo an ninh.}

\vocabentry{Affordability}
{/əˌfɔːrdəˈbɪləti/ (n)}
{The fact of being cheap enough that people can afford to buy it or pay for it.}
{Improving the affordability of green energy is essential for a fair transition.}
{Đây là một danh từ dùng để chỉ tính dễ tiếp cận về giá cả của năng lượng sạch}

\vocabentry{Standardization}
{/ˌstændərdəˈzeɪʃn/ (n)}
{The process of making things of the same type all have the same basic features.}
{International standardization of battery recycling processes is needed to protect the environment.}
{Đây là một danh từ dùng để chỉ sự áp dụng tiêu chuẩn chung cho các thiết bị và hệ thống năng lượng}

\vocabentry{Implementation}
{/ˌɪmplosámanˈteɪʃn/ (n)}
{The process of putting a decision or plan into effect; execution.}
{The slow implementation of carbon taxes has delayed the shift to green energy.}
{Đây là một danh từ dùng để chỉ sự thực hiện chính sách hoặc dự án năng lượng}

\vocabentry{Utilization}
{/ˌjuːtələˈzeɪʃn/ (n)}
{The action of making practical and effective use of something.}
{The utilization of waste heat from power plants can provide district heating for entire cities.}
{Đây là một danh từ dùng để chỉ sự tận dụng tối đa nguồn lực năng lượng hiện có}

\vocabentry{Compliance}
{/kəmˈplaɪəns/ (n)}
{The action or fact of complying with a wish or command.}
{Companies must ensure compliance with all safety regulations when building nuclear plants.}
{Đây là một danh từ dùng để chỉ sự tuân thủ quy định và tiêu chuẩn an toàn năng lượng}

\vocabentry{Feasible}
{/ˈfiːzəbl/ (adj)}
{Possible to do easily or conveniently.}
{It is now technically feasible to power small islands entirely with renewable energy.}
{Đây là một tính từ dùng để chỉ tính khả thi.}

\vocabentry{Drastic}
{/ˈdræstɪk/ (adj)}
{Likely to have a strong or far-reaching effect; radical and extreme.}
{Drastic measures are needed to reduce carbon emissions by fifty percent this decade.}
{Đây là một tính từ dùng để chỉ sự mạnh mẽ hoặc quyết liệt trong thay đổi chính sách năng lượng.}

\vocabentry{Fundamental}
{/ˌfʌndəˈmentl/ (adj)}
{Forming a necessary base or core; of central importance.}
{A fundamental shift in how we produce energy is necessary for a greener future.}
{Đây là một tính từ dùng để chỉ tính chất nền tảng.}

\vocabentry{Significant}
{/sɪɡˈnɪfɪkənt/ (adj)}
{Sufficiently great or important to be worthy of attention.}
{There has been a significant drop in the price of wind energy technology.}
{Đây là một tính từ dùng để chỉ sự đáng kể hoặc quan trọng.}

\vocabentry{Widespread}
{/ˌwaɪdˈspred/ (adj)}
{Found or distributed over a large area or number of people.}
{The widespread use of LED bulbs has greatly reduced energy consumption for lighting.}
{Đây là một tính từ dùng để chỉ sự phổ biến rộng rãi.}

\vocabentry{Paradigm}
{/ˈpærədaɪm/ (n)}
{A typical example or pattern of something; a model. (Đây là một danh từ dùng để chỉ một mô hình hoặc chuẩn mực (ví dụ: paradigm shift - chuyển đổi mô hình).).}
{We are seeing a paradigm shift from centralized to decentralized energy production.}
{}

\vocabentry{Technological}
{/ˌteknəˈlɑːdʒɪkl/ (adj)}
{Relating to or using technology.}
{Technological barriers still exist for long-duration energy storage.}
{Đây là một tính từ dùng để chỉ các yếu tố thuộc về công nghệ.}

\vocabentry{Environmental}
{/ɪnˌvaɪrənˈmentl/ (adj)}
{Relating to the natural world and the impact of human activity on its condition.}
{Every energy project must undergo a thorough environmental impact assessment.}
{Đây là một tính từ dùng để chỉ các yếu tố thuộc về môi trường.}

\vocabentry{Economic}
{/ˌekəˈnɑːmɪk/ (adj)}
{Relating to economics or the economy.}
{The economic viability of solar power has made it the cheapest source of electricity in many areas.}
{Đây là một tính từ dùng để chỉ tính kinh tế.}

\vocabentry{Potential}
{/pəˈtenʃl/ (n/adj)}
{Having or showing the capacity to become or develop into something in the future.}
{The country has enormous potential for the development of wave energy.}
{Đây là một danh từ/tính từ dùng để chỉ tiềm năng.}

\vocabentry{Mechanism}
{/ˈmekənɪzəm/ (n)}
{A system of parts working together in a machine; a process by which something takes place.}
{The mechanism of a heat pump allows it to provide both heating and cooling efficiently.}
{Đây là một danh từ dùng để chỉ cơ chế hoạt động.}

\vocabentry{Component}
{/kəmˈpoʊnənt/ (n)}
{A part or element of a larger whole.}
{Inverters are a key component of any solar power system.}
{Đây là một danh từ dùng để chỉ linh kiện hoặc bộ phận.}

\vocabentry{Stability}
{/stəˈbɪləti/ (n)}
{The state of being stable.}
{Grid stability is maintained by balancing electricity supply and demand at all times.}
{Trong năng lượng là sự ổn định của lưới điện.}

\vocabentry{Interconnection}
{/ˌɪntərkəˈnekʃn/ (n)}
{A mutual connection between two or more things.}
{The interconnection of regional grids helps to prevent blackouts.}
{Đây là một danh từ chỉ sự kết nối liên thông.}

\vocabentry{Connectivity}
{/ˌkɑːnekˈtɪvəti/ (n)}
{The state or extent of being connected or interconnected.}
{Smart appliances require internet connectivity to communicate with the smart meter.}
{Đây là một danh từ dùng để chễ khả năng kết nối của mạng lưới điện thông minh và thiết bị}

\vocabentry{Accessibility}
{/əkˌsesəˈbɪləti/ (n)}
{The quality of being able to be reached or entered.}
{Improving the accessibility of clean energy in developing nations is a global goal.}
{Đây là một danh từ dùng để chỉ khả năng tiếp cận năng lượng của người dân.}

\vocabentry{Availability}
{/əˌveɪləˈbɪləti/ (n)}
{The state of being otherwise unoccupied or at liberty to be used.}
{The availability of natural sunlight is the only limit to solar power production.}
{Đây là một danh từ dùng để chỉ sự sẵn có của nguồn năng lượng.}

\vocabentry{Reliability}
{/rɪˌlaɪəˈbɪləti/ (n)}
{The quality of being trustworthy or of performing consistently well.}
{Ensuring the reliability of the grid during extreme weather events is a major challenge.}
{Đây là một danh từ dùng để chỉ độ ổn định và tin cậy của hệ thống điện, không mất điện}

\vocabentry{Implement}
{/ˈɪmplɪment/ (v)}
{Put (a decision, plan, agreement, etc.) into effect.}
{The city plans to implement a new policy to ban fossil fuel heating in new homes.}
{Đây là một động từ dùng để chỉ việc thực hiện/áp dụng một giải pháp.}

\vocabentry{Maintenance}
{/ˈmeɪntənəns/ (n)}
{The process of maintaining or preserving someone or something.}
{Regular maintenance is required to keep wind turbines operating at peak efficiency.}
{Đây là một danh từ dùng để chỉ sự bảo trì thiết bị năng lượng.}

\vocabentry{Development}
{/dɪˈveləpmənt/ (n)}
{The process of developing or being developed.}
{The development of tidal energy technology is still in its early stages in many countries.}
{Đây là một danh từ dùng để chỉ sự phát triển/nghiên cứu mới.}

\vocabentry{Research}
{/rɪˈsɜːrtʃ/ (n/v)}
{The systematic investigation into and study of materials and sources.}
{Millions of dollars are spent on research into more efficient solar cells every year.}
{Đây là một danh từ/động từ dùng để chỉ hoạt động nghiên cứu.}

\vocabentry{Technology}
{/tekˈnɑːlədʒi/ (n)}
{The application of scientific knowledge for practical purposes.}
{New technology has made it possible to extract natural gas from deep underground.}
{Đây là một danh từ dùng để chỉ công nghệ.}

\vocabentry{Integration}
{/ˌɪntɪˈɡreɪʃn/ (n)}
{The action or process of integrating.}
{The integration of electric vehicles into the grid can help to balance demand.}
{Đây là một danh từ chỉ sự tích hợp các nguồn năng lượng sạch vào lưới điện.}

\vocabentry{Coordination}
{/koʊˌɔːrdɪˈneɪʃn/ (n)}
{The organization of the different elements of a complex body or activity.}
{Regional coordination is necessary to ensure a stable supply of electricity.}
{Đây là một danh từ chỉ sự phối hợp nhịp nhàng giữa các nhà máy điện.}

\vocabentry{Transformation}
{/ˌtrænsfərˈmeɪʃn/ (n)}
{A thorough or dramatic change in form or appearance.}
{We are witnessing a historic transformation of the global energy sector.}
{Đây là một danh từ chỉ sự biến đổi hoàn toàn cơ cấu năng lượng.}

\vocabentry{Revolution}
{/ˌrevəˈluːʃn/ (n)}
{A forcible overthrow of a government; in technology, a sudden and dramatic change.}
{The solar revolution has made clean energy accessible to millions of people.}
{Đây là một danh từ chỉ cuộc cách mạng năng lượng.}

\vocabentry{Advancement}
{/ədˈvænsmənt/ (n)}
{The process of promoting a cause or plan; a development or improvement.}
{Advancements in digital technology have enabled the creation of smart energy grids.}
{Đây là một danh từ chỉ sự tiến bộ.}

\vocabentry{Reduction}
{/rɪˈdʌkʃn/ (n)}
{The action or fact of making a specified thing smaller or less in amount. (Đây là một danh từ chỉ sự giảm thiểu (ví dụ: giảm khí thải).).}
{A significant reduction in methane leaks from gas pipelines is essential.}
{}

\vocabentry{Elimination}
{/ɪˌlɪmɪˈneɪʃn/ (n)}
{The complete removal or destruction of something. (Đây là một danh từ chỉ sự loại bỏ hoàn toàn (ví dụ: loại bỏ điện than).).}
{Many countries are aiming for the total elimination of coal power by 2040.}
{}

\vocabentry{Mitigation}
{/ˌmɪtɪˈɡeɪʃn/ (n)}
{The action of reducing the severity, seriousness, or painfulness of something.}
{Mitigation strategies in the energy sector are essential for a sustainable future.}
{Đây là một danh từ dùng để chỉ sự giảm nhẹ tác hại môi trường, đặc biệt là khí thải}

\vocabentry{Preservation}
{/ˌprezərˈveɪʃn/ (n)}
{The action of preserving something.}
{The preservation of local ecosystems is a major concern when building large dams.}
{Đây là một danh từ dùng để chỉ sự giữ gìn nguyên trạng môi trường khi khai thác năng lượng.}

\vocabentry{Protection}
{/prəˈtekʃn/ (n)}
{The action of protecting someone or something.}
{Stronger protection for the power grid against cyberattacks is necessary.}
{Đây là một danh từ chỉ sự bảo vệ môi trường hoặc hạ tầng năng lượng.}

\vocabentry{Restoration}
{/ˌrestəˈreɪʃn/ (n)}
{The action of returning something to a former owner, place, or condition.}
{The restoration of old coal mining sites can create new green spaces.}
{Đây là một danh từ chỉ sự phục hồi môi trường sau khai thác mỏ.}

\vocabentry{Adjustment}
{/əˈdʒʌstmənt/ (n)}
{A small alteration or movement made to achieve a desired fit or appearance.}
{Frequent adjustments to electricity production are needed to match fluctuating demand.}
{Đây là một danh từ chỉ sự điều chỉnh kỹ thuật.}

\vocabentry{Modification}
{/ˌmɑːdɪfɪˈkeɪʃn/ (n)}
{The action of modifying something.}
{Slight modifications to the blade design can make wind turbines much quieter.}
{Đây là một danh từ chỉ sự sửa đổi hoặc cải tiến thiết bị.}

\vocabentry{Alternating current}
{/ˌɔːltərneɪtɪŋ ˈkɜːrənt/ (n)}
{An electric current that reverses its direction many times a second at regular intervals.}
{Alternating current is the standard form of electricity used in homes and offices.}
{Đây là một danh từ dùng để chỉ dòng điện xoay chiều - AC.}

\vocabentry{Direct current}
{/dəˌrekt ˈkɜːrənt/ (n)}
{An electric current flowing in one direction only.}
{Batteries store energy as direct current, which must be converted for most appliances.}
{Đây là một danh từ dùng để chỉ dòng điện một chiều - DC.}

\vocabentry{Ampere}
{/ˈæmper/ (n)}
{A unit of electric current equal to a flow of one coulomb per second.}
{The fuse is designed to blow if the current exceeds fifteen amperes.}
{Đây là một danh từ dùng để chỉ đơn vị đo cường độ dòng điện - Ampe.}

\vocabentry{Watt}
{/wɑːt/ (n)}
{The SI unit of power, equivalent to one joule per second.}
{LED bulbs consume much fewer watts than old-fashioned light bulbs.}
{Đây là một danh từ dùng để chỉ đơn vị đo công suất - Oát.}

\vocabentry{Joule}
{/dʒuːl/ (n)}
{The SI unit of work or energy.}
{A joule is defined as the work done when a force of one newton moves an object one meter.}
{Đây là một danh từ dùng để chỉ đơn vị đo năng lượng - Jun.}

\vocabentry{Kilowatt-hour}
{/ˌkɪləwɑːt ˈaʊər/ (n)}
{A unit of energy equal to one kilowatt of power sustained for one hour.}
{Your electricity bill is based on the number of kilowatt-hours you use each month.}
{Đây là một danh từ dùng để chỉ số điện tiêu thụ - kWh.}

\vocabentry{Grid-connected}
{/ˈɡrɪd kəˈnektɪd/ (adj)}
{(Of a renewable energy system) linked to the local electrical power grid.}
{Most residential solar systems are grid-connected to allow for the sale of excess power.}
{Đây là một tính từ dùng để chỉ các hệ thống năng lượng được nối lưới điện.}

\vocabentry{Base load}
{/ˈbeɪs loʊd/ (n)}
{The minimum amount of electric power delivered or required over a given period at a constant rate.}
{Nuclear and coal power have traditionally provided the base load for the grid.}
{Đây là một cụm danh từ dùng để chỉ phụ tải cơ bản/điện nền.}

\vocabentry{Peak demand}
{/ˌpiːk dɪˈmænd/ (n)}
{The highest amount of electricity being used at any one time during a day.}
{Peak demand usually occurs in the early evening when people return home and turn on appliances.}
{Đây là một cụm danh từ dùng để chỉ nhu cầu cao điểm.}

\vocabentry{Blackout}
{/ˈblækaut/ (n)}
{A failure of electrical power supply.}
{The severe storm caused a major blackout that affected over a million people.}
{Đây là một danh từ dùng để chỉ sự mất điện diện rộng.}

\vocabentry{Brownout}
{/ˈbraʊnaʊt/ (n)}
{A reduction in or restriction on the availability of electrical power in a particular area.}
{During the heatwave, the utility company implemented brownouts to prevent a total grid failure.}
{Đây là một danh từ dùng để chỉ sự sụt áp điện.}

\vocabentry{Cogeneration}
{/ˌkoʊˌdʒenəˈreɪʃn/ (n)}
{The generation of electricity and useful heat jointly.}
{Cogeneration plants are much more efficient than traditional power plants because they use waste heat.}
{Đây là một danh từ dùng để chỉ quá trình đồng phát điện - nhiệt.}

\vocabentry{District heating}
{/ˌdɪstrɪkt ˈhiːtɪŋ/ (n)}
{A system for distributing heat generated in a centralized location through a system of insulated pipes.}
{District heating is a very efficient way to warm homes in cold northern cities.}
{Đây là một cụm danh từ dùng để chỉ hệ thống sưởi ấm khu vực.}

\vocabentry{Heat pump}
{/ˈhiːt pʌmp/ (n)}
{A device that transfers heat from a colder area to a hotter area using mechanical energy.}
{Installing a heat pump can be a very effective way to reduce carbon emissions from home heating.}
{Đây là một cụm danh từ dùng để chỉ máy bơm nhiệt.}

\vocabentry{Inverter}
{/ɪnˈvɜːrtər/ (n)}
{An electronic device or circuitry that changes direct current (DC) to alternating current (AC).}
{The inverter is a crucial part of a solar power system, allowing the energy to be used by home appliances.}
{Đây là một danh từ dùng để chỉ bộ biến tần điện.}

\vocabentry{Solar cell}
{/ˈsoʊlər sel/ (n)}
{A device converting solar radiation into electricity.}
{Modern solar cells are becoming much more efficient at capturing sunlight on cloudy days.}
{Đây là một cụm danh từ dùng để chỉ tế bào quang điện/pin mặt trời.}

\vocabentry{Concentrated solar}
{/ˌkɑːnsntreɪtɪd ˈsoʊlər/ (n)}
{A system that uses mirrors or lenses to focus a large area of sunlight onto a receiver.}
{Concentrated solar power plants can store heat to generate electricity even after the sun goes down.}
{Đây là một cụm danh từ dùng để chỉ năng lượng mặt trời tập trung.}

\vocabentry{Passive solar}
{/ˌpæsɪv ˈsoʊlər/ (adj)}
{(Of a building) designed to use the sun's energy for heating and cooling without active systems.}
{Passive solar design includes large south-facing windows and materials that store heat.}
{Đây là một tính từ dùng để chỉ thiết kế năng lượng mặt trời thụ động.}

\vocabentry{Wind vane}
{/ˈwɪnd veɪn/ (n)}
{A device that measures wind direction.}
{The wind vane on top of the turbine allows it to turn and face the wind for maximum efficiency.}
{Đây là một danh từ dùng để chỉ phong kế/cái chỉ hướng gió.}

\vocabentry{Nacelle}
{/nəˈsel/ (n)}
{The housing that contains the generating components in a wind turbine.}
{The nacelle contains the gearbox, generator, and control systems of the wind turbine.}
{Đây là một danh từ chuyên ngành dùng để chỉ vỏ động cơ của tua-bin gió.}

\vocabentry{Rotor}
{/ˈroʊtər/ (n)}
{The rotating part of a turbine, electric motor, or other machine.}
{The massive rotor of the wind turbine can be over one hundred meters in diameter.}
{Đây là một danh từ dùng để chỉ bộ phận quay/rô-to.}

\vocabentry{Dam}
{/dæm/ (n)}
{A barrier constructed to hold back water and raise its level.}
{Building a new dam can provide clean energy but may also displace local communities.}
{Đây là một danh từ dùng để chỉ đập thủy điện.}

\vocabentry{Reservoir}
{/ˈrezərvwɑːr/ (n)}
{A large natural or artificial lake used as a source of water supply.}
{The reservoir behind the dam is also used for recreation and irrigation.}
{Đây là một danh từ dùng để chỉ hồ chứa nước cho thủy điện.}

\vocabentry{Penstock}
{/ˈpenstɑːk/ (n)}
{A sluice or gate or intake structure that controls water flow to a turbine.}
{Water flows from the reservoir through the penstock to spin the turbines at high speed.}
{Đây là một danh từ chuyên ngành dùng để chỉ ống dẫn nước áp lực trong thủy điện.}

\vocabentry{Geothermal well}
{/ˌdʒiːoʊˈθɜːrml wel/ (n)}
{A hole drilled into the earth to tap into geothermal energy.}
{Geothermal wells can reach depths of several kilometers to find hot water or steam.}
{Đây là một cụm danh từ dùng để chỉ giếng địa nhiệt.}

\vocabentry{Dry steam}
{/ˌdraɪ ˈstiːm/ (n)}
{Steam that contains no water droplets.}
{Dry steam power plants are the simplest and oldest type of geothermal technology.}
{Trong địa nhiệt, đây là hơi khô dùng để chạy tua-bin trực tiếp.}

\vocabentry{Flash steam}
{/ˌflæʃ ˈstiːm/ (n)}
{Steam produced when high-pressure hot water is released into a lower-pressure tank.}
{Most modern geothermal plants use flash steam technology to generate electricity.}
{Đây là một loại công nghệ địa nhiệt phổ biến.}

\vocabentry{Binary cycle}
{/ˌbaɪnəri ˈsaɪkl/ (n)}
{A geothermal power plant that uses a second fluid with a lower boiling point than water.}
{Binary cycle plants can generate electricity from lower-temperature geothermal resources.}
{Đây là công nghệ địa nhiệt dùng vòng tuần hoàn kép.}

\vocabentry{Anaerobic}
{/ˌænerˈoʊbɪk/ (adj)}
{Relating to or requiring an absence of free oxygen.}
{Anaerobic digestion of food waste produces methane that can be used for cooking or heating.}
{Trong năng lượng sinh khối, đây là quá trình phân hủy kỵ khí để tạo khí biogas.}

\vocabentry{Gasification}
{/ˌɡæsɪfɪˈkeɪʃn/ (n)}
{A process that converts organic or fossil-based carbonaceous materials into carbon monoxide and hydrogen.}
{Biomass gasification is a very efficient way to produce clean-burning fuel from wood waste.}
{Đây là một danh từ dùng để chỉ quá trình khí hóa nhiên liệu.}

\vocabentry{Pyrolysis}
{/paɪˈrɑːləsɪs/ (n)}
{Decomposition brought about by high temperatures in the absence of oxygen.}
{Pyrolysis can convert agricultural waste into a liquid fuel similar to crude oil.}
{Đây là một danh từ dùng để chỉ quá trình nhiệt phân tạo ra dầu sinh học.}

\vocabentry{Ethanol}
{/ˈeθənɔːl/ (n)}
{A colorless volatile flammable liquid which is produced by the natural fermentation of sugars.}
{Most gasoline sold in the US contains ten percent ethanol made from corn.}
{Đây là một danh từ dùng để chỉ cồn nhiên liệu.}

\vocabentry{Biodiesel}
{/ˌbaɪoʊˈdiːzl/ (n)}
{A biofuel intended as a substitute for diesel, produced from vegetable oil or animal fat.}
{Biodiesel can be used in most modern diesel engines without any modifications.}
{Đây là một danh từ dùng để chỉ dầu đi-ê-zen sinh học.}

\vocabentry{Radioactive}
{/ˌreɪdioʊˈæktɪv/ (adj)}
{Emitting or relating to the emission of ionizing radiation or particles.}
{Radioactive waste from nuclear plants must be stored safely for thousands of years.}
{Đây là một tính từ dùng để chỉ tính phóng xạ của nhiên liệu hạt nhân.}

\vocabentry{Enrichment}
{/ɪnˈrɪtʃmənt/ (n)}
{The action of improving or enhancing the quality or value of something.}
{Uranium enrichment is a complex process required to produce fuel for nuclear reactors.}
{Trong hạt nhân, đây là sự làm giàu uranium.}

\vocabentry{Meltdown}
{/ˈmeltdown/ (n)}
{A severe nuclear reactor accident that results in core damage from overheating.}
{The Fukushima disaster was caused by a partial nuclear meltdown following a tsunami.}
{Đây là một danh từ dùng để chỉ sự nóng chảy lõi lò phản ứng hạt nhân.}

\vocabentry{Isotope}
{/ˈaɪsətoʊp/ (n)}
{Each of two or more forms of the same element that contain equal numbers of protons but different numbers of neutrons.}
{Uranium-235 is the isotope used as fuel in most nuclear power plants.}
{Đây là một danh từ dùng để chỉ đồng vị hạt nhân.}

\vocabentry{Half-life}
{/ˈhæf laɪf/ (n)}
{The time taken for the radioactivity of a specified isotope to fall to half its original value.}
{Some radioactive waste has a half-life of only a few days, while others last for centuries.}
{Đây là một danh từ dùng để chỉ chu kỳ bán rã của chất phóng xạ.}

\vocabentry{Ocean thermal}
{/ˌoʊʃn ˈθɜːrml/ (n)}
{Energy generated from the temperature difference between deep cold ocean water and warm surface waters.}
{Ocean thermal energy conversion is a promising source of renewable power for tropical island nations.}
{Đây là năng lượng nhiệt đại dương - OTEC.}

\vocabentry{Wave energy}
{/ˈweɪv ˌenərdʒi/ (n)}
{The transport and capture of energy by ocean surface waves.}
{Wave energy converters use the movement of the ocean surface to generate electricity.}
{Đây là một cụm danh từ dùng để chỉ năng lượng sóng biển.}

\vocabentry{Osmotic power}
{/ɑːzˈmɑːtɪk ˈpaʊər/ (n)}
{Energy available from the difference in the salt concentration between seawater and river water.}
{The first osmotic power plant was built in Norway to test the feasibility of the technology.}
{Đây là năng lượng thẩm thấu.}

\vocabentry{Piezoelectric}
{/ˌpiːeɪzoʊɪˈlektrɪk/ (adj)}
{Relating to electric polarization in a substance resulting from the application of mechanical stress.}
{Piezoelectric materials could be used in floors to generate electricity from people walking.}
{Đây là một tính từ dùng để chỉ hiệu ứng áp điện, tạo điện từ áp lực.}

\vocabentry{Thermoelectric}
{/ˌθɜːrmoʊɪˈlektrɪk/ (adj)}
{Relating to the direct conversion of temperature differences to electric voltage.}
{Thermoelectric generators can capture waste heat from car exhausts to power electronics.}
{Đây là một tính từ dùng để chỉ hiệu ứng nhiệt điện.}

\vocabentry{Superconductivity}
{/ˌsuːpərˌkɑːndʌkˈtɪvəti/ (n)}
{The property of zero electrical resistance in some substances at very low temperatures.}
{Superconductivity could allow for the transmission of electricity over long distances with zero loss.}
{Đây là một danh từ dùng để chỉ hiện tượng siêu dẫn.}

\vocabentry{Carbon neutral}
{/ˌkɑːrbən ˈnuːtrəl/ (adj)}
{Making no net release of carbon dioxide to the atmosphere.}
{Many tech giants aim to become carbon neutral by investing in large wind and solar projects.}
{Đây là một tính từ dùng để chỉ trạng thái trung hòa carbon.}

\vocabentry{Carbon sink}
{/ˈkɑːrbən sɪŋk/ (n)}
{A forest, ocean, or other natural environment viewed in terms of its ability to absorb carbon dioxide.}
{Protecting existing forests is important because they act as vital carbon sinks.}
{Đây là một cụm danh từ dùng để chỉ bể chứa carbon tự nhiên.}

\vocabentry{Sequestration}
{/ˌsiːkwəˈstreɪʃn/ (n)}
{The action of taking legal possession of assets; in energy, the long-term storage of carbon.}
{Biological carbon sequestration involves planting millions of trees to absorb CO2.}
{Đây là danh từ chỉ sự cô lập hoặc lưu trữ carbon lâu dài.}

\vocabentry{Vulnerability}
{/ˌvʌlnərəˈbɪləti/ (n)}
{The quality or state of being exposed to the possibility of being attacked or harmed.}
{The vulnerability of centralized power grids to natural disasters is a major concern.}
{Sự dễ bị tổn thương của hệ thống năng lượng.}

